\begin{abcliste}
\abc A bulb does not use up the current: all the charge that flows into it flows out of it again. What the bulb turns into light and heat is energy, at the rate $P = \Delta V\cdot I$. So ammeter 2, behind A, reads the same as ammeter 1:
\begin{align}
 I_2 &= I_1 = \result{\IaP{0}{2}}
\end{align}
At the junction the current splits. B and C are in parallel, so they have the same voltage; they are identical, so they let the same current through. The current divides equally:
\begin{align}
 I_3 = I_4 &= \IbF = \frac{\Ia}{2} = \result{\IbP{0}{2}}
\end{align}
Behind the pair the two currents join again (junction rule):
\begin{align}
 I_5 &= I_3 + I_4 = I_1 = \result{\IaP{0}{2}}
\end{align}

\abc The same current flows through D as through A. Identical bulbs with the same current have the same voltage across them and so the same power $P = \Delta V\cdot I$: D is \result{\text{as bright as A}}. (A typical mistake is to think that D gets less current because A, B and C have used some of it up.) B and C are less bright than A and D: each carries only half the current, at a smaller voltage.
\end{abcliste}
